\documentclass[12pt,a4paper]{article}
\usepackage[utf8]{inputenc}
\usepackage[T2A]{fontenc}
\usepackage[russian]{babel}
\usepackage{amsmath,amssymb}
\usepackage{array}
\usepackage{longtable}
\usepackage{booktabs}
\usepackage{geometry}
\usepackage{hyperref}
\geometry{margin=2.3cm}

\title{Аудит гипотезы общей тени для \texttt{TOE.pdf} и \texttt{UGSM}:\
новый документ по численной проверке, результатам и интерпретации}
\author{}
\date{}

\begin{document}

\maketitle

\begin{abstract}
В этом документе фиксируется именно тот аудит, который был проведён для проверки гипотезы, что \texttt{TOE.pdf} и \texttt{ugsm\_dynamics\_audit-3.pdf} могут быть двумя тенями одной более общей спектральной сущности. В отличие от общего сравнительного анализа, здесь собраны: постановка гипотезы, использованная математическая редукция, конкретный численный эксперимент, численные результаты с цифрами, их интерпретация и итоговый вердикт о том, насколько правдоподобна гипотеза общего источника. Статус вывода --- исследовательский аудит с численной поддержкой, но не окончательное доказательство.
\end{abstract}

\section{Что именно проверялось}

Проверялся следующий тезис:
\begin{quote}
существует ли один более общий спектральный операторный объект, для которого корреляционный блок TOE и спектральный функционал UGSM являются двумя разными представлениями --- тепловым и моментным.
\end{quote}

В рабочей форме гипотеза выглядела так:
\begin{equation}
\hat C_\sigma \sim e^{-\sigma D_M^2/2},
\end{equation}
где \(D_M^2\) --- оператор Лаплас/Дирак-типа на компактном многообразии вакуумного сектора UGSM, а именно на \(M=S^3\times S^1\).

Если такая редукция допустима, то:
\begin{itemize}
    \item TOE читает тот же спектр в форме теплового корреляционного ядра;
    \item UGSM читает тот же спектр в форме спектрального функционала
    \[
    S_{spec}(\Lambda)=\mathrm{Tr}\,f(D_M^2/\Lambda^2);
    \]
    \item обе теории оказываются не независимыми конструкциями, а двумя тенями одного более общего спектрального уровня.
\end{itemize}

\section{Какая проверка была проведена}

Для первичной проверки был выбран минимальный и контролируемый тест на \(S^3\times S^1\): скалярный тепловой след лапласиана. Это не доказывает полную эквивалентность TOE и UGSM, но позволяет проверить ключевой признак: воспроизводит ли один и тот же тепловой объект те геометрические коэффициенты, которые естественно поддерживают UGSM-тип спектрального описания.

Использовалась следующая схема.

\subsection{Геометрия и спектр}

Был взят продукт
\[
M=S^3(R_3)\times S^1(R_1)
\]
с радиусами
\[
R_3=1, \qquad R_1=1.
\]

Для скалярного лапласиана использовались собственные значения и кратности:
\begin{align}
\lambda_{n,m} &= \frac{n(n+2)}{R_3^2}+\frac{m^2}{R_1^2},\\
 g_n &= (n+1)^2,
\end{align}
где \(n\ge 0\), \(m\in\mathbb{Z}\).

\subsection{Тепловой след}

Вычислялся тепловой след
\begin{equation}
K(t)=\sum_{n,m} g_n e^{-t\lambda_{n,m}}.
\end{equation}

Затем рассматривалась стандартная безразмерная комбинация
\begin{equation}
\widetilde K(t)=(4\pi t)^2 K(t),
\end{equation}
поскольку размерность многообразия равна четырём. Именно такая нормировка позволяет напрямую читать первые коэффициенты асимптотики теплового следа.

\subsection{Численная реализация}

Численный аудит был выполнен с усечением спектра
\[
n_{\max}=80, \qquad m_{\max}=80.
\]

Фит по малым \(t\) проводился на окне
\[
t\in\{0.03,0.035,0.04,0.045,0.05,0.055,0.06\}.
\]

Использовалась квадратичная аппроксимация
\begin{equation}
\widetilde K(t)\approx a_0+a_2 t+a_4 t^2,
\end{equation}
где коэффициенты \(a_0,a_2,a_4\) извлекались методом наименьших квадратов.

\section{Что было найдено численно}

Из проведённого расчёта получены следующие значения.

\begin{longtable}{p{6.0cm}p{4.2cm}p{4.2cm}}
\toprule
Параметр & Численное значение & Теоретическое значение / комментарий \\
\midrule
\endhead
Многообразие & \(S^3\times S^1\) & минимальный вакуумный фон UGSM \\
Оператор & скалярный лапласиан & минимальный тестовый выбор \\
Радиусы & \(R_3=1,\ R_1=1\) & единичная нормировка \\
Усечение спектра & \(n_{\max}=80,\ m_{\max}=80\) & достаточно для устойчивого малого окна \\
Окно фита & \(0.03\) -- \(0.06\) & область малого \(t\) \\
\(a_0^{fit}\) & 124.0268873911 & ожидается \(4\pi^3=124.0251067212\) \\
\(a_2^{fit}\) & 123.8993672201 & ожидается \(4\pi^3=124.0251067212\) \\
\(a_4^{fit}\) & 64.8681519650 & следующий кривизновый вклад \\
\(a_2^{fit}/(4\pi)\) & 9.8595983695 & ожидается \(\pi^2=9.8696044011\) \\
Систола окружности & 6.2831853072 & равна \(2\pi\) при \(R_1=1\) \\
\(\pi\)-proxy из систолы & 3.1415926536 & равна \(\pi\) \\
\bottomrule
\end{longtable}

\section{Точность совпадений}

Для ключевых коэффициентов были вычислены абсолютные и относительные отклонения.

\begin{longtable}{p{4.8cm}p{3.4cm}p{3.8cm}p{3.4cm}}
\toprule
Величина & Численное значение & Ожидаемое значение & Относительная ошибка \\
\midrule
\endhead
\(a_0\) & 124.0268873911 & 124.0251067212 & \(1.44\times 10^{-5}\) \\
\(a_2\) & 123.8993672201 & 124.0251067212 & \(1.01\times 10^{-3}\) \\
\(a_2/(4\pi)\) & 9.8595983695 & 9.8696044011 & \(1.01\times 10^{-3}\) \\
\bottomrule
\end{longtable}

Это даёт две важные численные констатации:
\begin{enumerate}
    \item ведущий коэффициент \(a_0\) воспроизводится практически идеально уже в простейшей схеме;
    \item следующий коэффициент \(a_2\) и его нормированная версия \(a_2/(4\pi)\) восстанавливаются с ошибкой порядка \(10^{-3}\), что для такого усечённого и минимального теста является очень сильным совпадением.
\end{enumerate}

\section{Как интерпретируется найденное}

\subsection{Первый уровень интерпретации: общая геометрическая тень}

Численно найдено, что один и тот же тепловой объект на \(S^3\times S^1\) естественно порождает масштабы:
\[
4\pi^3, \qquad \pi^2, \qquad \pi.
\]

Это важно потому, что именно такие \(\pi\)-структуры выглядят не случайными константами, а геометрическими отпечатками компактного спектрального фона. Следовательно, значимая часть UGSM-тип феноменологии действительно может интерпретироваться как чтение тех же геометрических данных, которые в TOE естественно появляются через тепловое корреляционное ядро.

\subsection{Второй уровень интерпретации: две репрезентации одного спектра}

Если принять рабочую редукцию
\[
\hat C_\sigma \sim e^{-\sigma D_M^2/2},
\]
то:
\begin{itemize}
    \item TOE читает зависимость по параметру теплового времени \(\sigma\);
    \item UGSM читает моменты того же спектра через функцию \(f\) в спектральном действии;
    \item различие между теориями переносится с уровня ``два разных закона'' на уровень ``два способа читать один и тот же операторный контент''.
\end{itemize}

Именно это и понимается здесь под выражением ``общая тень'': не буквенная идентичность формул, а единый спектральный источник, видимый в разных функциональных проекциях.

\subsection{Третий уровень интерпретации: что означает линейный pi-сектор}

В численном аудите также естественно появляется систолическая величина окружности:
\[
\operatorname{systole}(S^1)=2\pi,
\]
откуда сразу следует простой \(\pi\)-proxy:
\[
\frac{\operatorname{systole}(S^1)}{2}=\pi.
\]

Это ещё не строгий вывод линейного \(\pi\)-вклада в полной физической теории, но это сильный индикатор того, что линейный сектор может иметь не произвольную фитовую природу, а топологическую или голономную интерпретацию. В языке моста это означает: \(\pi\)-вклад выглядит не как случайная подгонка, а как кандидат на геометрически защищённый след компактного цикла.

\section{Что именно этот аудит подтверждает}

Проведённый аудит подтверждает следующие более узкие, но важные пункты:
\begin{enumerate}
    \item гипотеза общего спектрального источника математически осмысленна;
    \item минимальный тепловой объект на \(S^3\times S^1\) действительно восстанавливает ключевые геометрические коэффициенты с высокой точностью;
    \item шкалы \(4\pi^3\), \(\pi^2\) и \(\pi\) могут быть поняты как части одной геометрической структуры, а не как независимые числовые совпадения;
    \item поэтому тезис ``TOE и UGSM --- две тени одной большой сущности'' получает не только философскую, но и численно поддержанную рабочую форму.
\end{enumerate}

\section{Что именно этот аудит пока не доказывает}

При всей силе найденных совпадений этот документ не доказывает ещё несколько критических вещей:
\begin{enumerate}
    \item не доказано, что полный корреляционный оператор TOE строго редуцируется именно к \(e^{-\sigma D_M^2/2}\) во всех релевантных секторах;
    \item не доказано, что полный функционал UGSM целиком восстанавливается из одной и той же меры Лапласова типа без внешних феноменологических ручек;
    \item не доказано, что линейный \(\pi\)-сектор строго выводится как голономный или систолический инвариант полной конструкции;
    \item не доказано, что все наблюдаемые UGSM и все ключевые блоки TOE являются лишь разными проекциями одного и того же оператора без дополнительных допущений.
\end{enumerate}

Поэтому язык ``общая тень'' в этом документе означает сильную рабочую гипотезу, поддержанную математическим мостом и численным аудитом, но пока не теорему об эквивалентности.

\section{Итоговый вердикт}

На основании проведённого аудита наиболее аккуратный итоговый вердикт таков:

\begin{quote}
С высокой степенью правдоподобия \texttt{TOE.pdf} и \texttt{ugsm\_dynamics\_audit-3.pdf} можно рассматривать как две разные тени одной более общей спектральной сущности. Этот вывод поддерживается не только концептуальным мостом, но и конкретной численной проверкой на \(S^3\times S^1\), где один и тот же тепловой след воспроизводит ключевые геометрические масштабы \(4\pi^3\), \(\pi^2\) и естественный \(\pi\)-proxy.
\end{quote}

Однако строгое утверждение о том, что обе теории действительно являются редукциями одной и той же большой сущности, пока делать рано. На текущем этапе корректнее говорить так:
\begin{itemize}
    \item гипотеза общей сущности стала математически конкретной;
    \item она получила сильную численную поддержку;
    \item она ещё требует строгого доказательного цикла для перехода от ``правдоподобного моста'' к ``доказанной редукции''.
\end{itemize}

\section{Практический вывод для следующего цикла}

Если продолжать этот проект дальше, следующий сильный шаг должен состоять в следующем:
\begin{itemize}
    \item заменить скалярный лапласиан на более физически релевантный оператор Дирака-типа;
    \item проверить устойчивость коэффициентов при варьировании радиусов \(R_3\) и \(R_1\);
    \item проверить, можно ли восстановить UGSM-подобные комбинации через более общий выбор тестовой функции \(f\);
    \item вынести линейный \(\pi\)-сектор из статуса proxy в статус строгого геометрического вывода.
\end{itemize}

Именно тогда вопрос ``являются ли обе теории тенью одной большой сущности?'' сможет перейти из режима сильной рабочей гипотезы в режим существенно более строгого утверждения.

\end{document}